\documentclass[10pt]{extarticle}
\usepackage[a4paper,margin=1.5cm]{geometry}
\usepackage[T1]{fontenc}
\usepackage{amsmath,amssymb}
\usepackage{graphicx}
\usepackage[export]{adjustbox}
\usepackage{multicol}
\usepackage{enumitem}
\usepackage{booktabs,array,tabularx}
\setlength{\columnsep}{0.7cm}
\setlength{\parindent}{0pt}
\setlength{\parskip}{3pt}
\pagestyle{empty}
\makeatletter
% figure registry (figdefs.tex is generated by the extract step)
\newcommand{\qfigdef}[4]{\@namedef{qf@file@#1}{#2}\@namedef{qf@w@#1}{#3}\@namedef{qf@wide@#1}{#4}}
\newcommand{\qf@check}[1]{\@ifundefined{qf@file@#1}{\PackageError{score2}{Unknown figure id #1}{}}{}}
\newcommand{\qf@img}[4][]{\includegraphics[width=#2\dimexpr\@nameuse{qf@w@#4}pt\relax,max width=#3,max height=0.45\textheight,#1]{\@nameuse{qf@file@#4}}}
% Full-width questions: finalize wraps every question that contains a wide figure in \qfullbegin...\qfullend,
% so the WHOLE question (stem, figures, options) is set across the page in reading order, then the two columns
% resume. The question's head - its number, stem and first figure (\qfullsplit marks the end of the head) -
% is kept on one page, so a number is never left at a page foot with its figure overleaf; the rest may break
% across pages normally, so no big gaps are left. Without a \qfullsplit, the whole question is kept together
% if it is under 45% of a page. Never write these yourself.
\newif\ifqf@full
\newif\ifqf@boxed
\newsavebox{\qf@box}
\newcommand{\qf@open}{\global\qf@fulltrue\par\global\qf@boxedtrue
  \setbox\qf@box\vbox\bgroup\hsize\textwidth\linewidth\textwidth}
\newcommand{\qf@emit}[1]{\par\egroup\global\qf@boxedfalse
  \ifdim\dimexpr\ht\qf@box+\dp\qf@box\relax<#1\textheight\noindent\usebox{\qf@box}\par\else\unvbox\qf@box\fi}
\newcommand{\qfullsplit}{\ifqf@boxed\qf@emit{0.85}\fi}
\newcommand{\qf@close}{\ifqf@boxed\qf@emit{0.45}\fi\global\qf@fullfalse}
\newcommand{\qfullbegin}{\end{multicols}\qf@open}
\newcommand{\qfullend}{\qf@close\begin{multicols}{2}}
\newcommand{\qfullstart}{\qf@open}
\newcommand{\qfulldone}{\qf@close}
% \qfig[scale]{id}: block figure, centred, directly where it is written (size set by the tool).
% Fallback if a wide figure is not inside a full-width question: break out of the columns in place - no floats.
\newcommand{\qfig}[2][1]{\qf@check{#2}%
  \ifqf@full\qf@blk{#1}{#2}\else
  \ifnum\@nameuse{qf@wide@#2}=1 \expandafter\@firstoftwo\else\expandafter\@secondoftwo\fi
  {\end{multicols}\begin{center}\qf@img{#1}{\textwidth}{#2}\end{center}\begin{multicols}{2}}%
  {\qf@blk{#1}{#2}}\fi}
\newcommand{\qf@blk}[2]{\par\smallskip{\centering\qf@img{#1}{\linewidth}{#2}\par}\smallskip}
% \qfigi[scale]{id}: inline figure (inside option lists / tables / minipages); never breaks columns.
\newcommand{\qfigi}[2][1]{\qf@check{#2}\qf@img[valign=c]{#1}{\linewidth}{#2}}
\makeatother
% structure macros - lint reads these to check completeness and figure placement
\newcommand{\papersec}[1]{\par\bigskip{\centering\Large\bfseries PAPER-#1\par}\medskip}
\newcommand{\partsec}[2]{\par\medskip\noindent{\large\bfseries PART-#1: #2}\par\smallskip}
\newcommand{\sect}[1]{\par\smallskip\noindent\textbf{#1}\par}
\newcommand{\qn}[1]{\par\smallskip\noindent\textbf{#1.}~}
\newcommand{\paraq}[2]{\par\smallskip\noindent\textbf{Paragraph for Questions #1 and #2}\par}
\newenvironment{qoptions}{\begin{itemize}[leftmargin=1.8em,itemsep=1pt,topsep=2pt,label={}]}{\end{itemize}}
\newcommand{\opt}[2]{\item[(#1)] #2}
% figure registry (generated by score2_tools.py finalize)
\qfigdef{p1_phy_s1_q01_1}{TOAS - Test 12 (10 Jul 2022)-img/p1_phy_s1_q01_1.png}{77.9}{0}
\qfigdef{p1_phy_s1_q02_1}{TOAS - Test 12 (10 Jul 2022)-img/p1_phy_s1_q02_1.png}{92.0}{0}
\qfigdef{p1_phy_s1_q03_1}{TOAS - Test 12 (10 Jul 2022)-img/p1_phy_s1_q03_1.png}{67.8}{0}
\qfigdef{p1_phy_s1_q04_1}{TOAS - Test 12 (10 Jul 2022)-img/p1_phy_s1_q04_1.png}{66.6}{0}
\qfigdef{p1_phy_s1_q05_1}{TOAS - Test 12 (10 Jul 2022)-img/p1_phy_s1_q05_1.png}{66.3}{0}
\qfigdef{p1_phy_s1_q06_1}{TOAS - Test 12 (10 Jul 2022)-img/p1_phy_s1_q06_1.png}{210.6}{0}
\qfigdef{p1_phy_s1_q08_1}{TOAS - Test 12 (10 Jul 2022)-img/p1_phy_s1_q08_1.png}{112.5}{0}
\qfigdef{p1_phy_s1_q09_1}{TOAS - Test 12 (10 Jul 2022)-img/p1_phy_s1_q09_1.png}{98.4}{0}
\qfigdef{p1_phy_s1_q10_1}{TOAS - Test 12 (10 Jul 2022)-img/p1_phy_s1_q10_1.png}{105.8}{0}
\qfigdef{p1_phy_s2_q01_1}{TOAS - Test 12 (10 Jul 2022)-img/p1_phy_s2_q01_1.png}{160.4}{0}
\qfigdef{p1_phy_s2_q03_1}{TOAS - Test 12 (10 Jul 2022)-img/p1_phy_s2_q03_1.png}{103.8}{0}
\qfigdef{p1_phy_s2_q05_1}{TOAS - Test 12 (10 Jul 2022)-img/p1_phy_s2_q05_1.png}{116.8}{0}
\qfigdef{p1_phy_s2_q06_1}{TOAS - Test 12 (10 Jul 2022)-img/p1_phy_s2_q06_1.png}{67.5}{0}
\qfigdef{p1_phy_s2_q08_1}{TOAS - Test 12 (10 Jul 2022)-img/p1_phy_s2_q08_1.png}{88.2}{0}
\qfigdef{p1_chem_s1_q04_1}{TOAS - Test 12 (10 Jul 2022)-img/p1_chem_s1_q04_1.png}{81.5}{0}
\qfigdef{p1_chem_s1_q05_A}{TOAS - Test 12 (10 Jul 2022)-img/p1_chem_s1_q05_A.png}{134.2}{0}
\qfigdef{p1_chem_s1_q05_A2}{TOAS - Test 12 (10 Jul 2022)-img/p1_chem_s1_q05_A2.png}{98.5}{0}
\qfigdef{p1_chem_s1_q05_B}{TOAS - Test 12 (10 Jul 2022)-img/p1_chem_s1_q05_B.png}{94.2}{0}
\qfigdef{p1_chem_s1_q05_C}{TOAS - Test 12 (10 Jul 2022)-img/p1_chem_s1_q05_C.png}{93.7}{0}
\qfigdef{p1_chem_s1_q06_1}{TOAS - Test 12 (10 Jul 2022)-img/p1_chem_s1_q06_1.png}{133.4}{0}
\qfigdef{p1_chem_s1_q07_A}{TOAS - Test 12 (10 Jul 2022)-img/p1_chem_s1_q07_A.png}{105.8}{0}
\qfigdef{p1_chem_s1_q07_B}{TOAS - Test 12 (10 Jul 2022)-img/p1_chem_s1_q07_B.png}{129.4}{0}
\qfigdef{p1_chem_s1_q08_1}{TOAS - Test 12 (10 Jul 2022)-img/p1_chem_s1_q08_1.png}{164.1}{0}
\qfigdef{p1_chem_s1_q08_A}{TOAS - Test 12 (10 Jul 2022)-img/p1_chem_s1_q08_A.png}{49.4}{0}
\qfigdef{p1_chem_s1_q08_B}{TOAS - Test 12 (10 Jul 2022)-img/p1_chem_s1_q08_B.png}{53.6}{0}
\qfigdef{p1_chem_s1_q08_C}{TOAS - Test 12 (10 Jul 2022)-img/p1_chem_s1_q08_C.png}{32.0}{0}
\qfigdef{p1_chem_s1_q08_D}{TOAS - Test 12 (10 Jul 2022)-img/p1_chem_s1_q08_D.png}{53.6}{0}
\qfigdef{p1_chem_s2_q07_1}{TOAS - Test 12 (10 Jul 2022)-img/p1_chem_s2_q07_1.png}{113.7}{0}
\qfigdef{p1_chem_s2_q08_2}{TOAS - Test 12 (10 Jul 2022)-img/p1_chem_s2_q08_2.png}{223.9}{0}
\qfigdef{p2_phy_s1_q01_1}{TOAS - Test 12 (10 Jul 2022)-img/p2_phy_s1_q01_1.png}{76.5}{0}
\qfigdef{p2_phy_s1_q03_1}{TOAS - Test 12 (10 Jul 2022)-img/p2_phy_s1_q03_1.png}{90.6}{0}
\qfigdef{p2_phy_s1_q04_1}{TOAS - Test 12 (10 Jul 2022)-img/p2_phy_s1_q04_1.png}{135.2}{0}
\qfigdef{p2_phy_s1_q06_1}{TOAS - Test 12 (10 Jul 2022)-img/p2_phy_s1_q06_1.png}{109.0}{0}
\qfigdef{p2_phy_s2_para7_8_1}{TOAS - Test 12 (10 Jul 2022)-img/p2_phy_s2_para7_8_1.png}{268.7}{1}
\qfigdef{p2_phy_s2_para9_10_1}{TOAS - Test 12 (10 Jul 2022)-img/p2_phy_s2_para9_10_1.png}{97.0}{0}
\qfigdef{p2_phy_s2_para11_12_1}{TOAS - Test 12 (10 Jul 2022)-img/p2_phy_s2_para11_12_1.png}{98.5}{0}
\qfigdef{p2_phy_s3_q01_1}{TOAS - Test 12 (10 Jul 2022)-img/p2_phy_s3_q01_1.png}{65.9}{0}
\qfigdef{p2_phy_s3_q02_1}{TOAS - Test 12 (10 Jul 2022)-img/p2_phy_s3_q02_1.png}{112.0}{0}
\qfigdef{p2_phy_s3_q03_1}{TOAS - Test 12 (10 Jul 2022)-img/p2_phy_s3_q03_1.png}{129.8}{0}
\qfigdef{p2_phy_s3_q04_1}{TOAS - Test 12 (10 Jul 2022)-img/p2_phy_s3_q04_1.png}{126.4}{0}
\qfigdef{p2_phy_s3_q05_1}{TOAS - Test 12 (10 Jul 2022)-img/p2_phy_s3_q05_1.png}{58.3}{0}
\qfigdef{p2_chem_s1_q04_1}{TOAS - Test 12 (10 Jul 2022)-img/p2_chem_s1_q04_1.png}{278.4}{1}
\qfigdef{p2_chem_s1_q05_1}{TOAS - Test 12 (10 Jul 2022)-img/p2_chem_s1_q05_1.png}{50.0}{0}
\qfigdef{p2_chem_s1_q05_A}{TOAS - Test 12 (10 Jul 2022)-img/p2_chem_s1_q05_A.png}{75.2}{0}
\qfigdef{p2_chem_s1_q05_B}{TOAS - Test 12 (10 Jul 2022)-img/p2_chem_s1_q05_B.png}{143.5}{0}
\qfigdef{p2_chem_s1_q05_C}{TOAS - Test 12 (10 Jul 2022)-img/p2_chem_s1_q05_C.png}{175.6}{0}
\qfigdef{p2_chem_s1_q05_D}{TOAS - Test 12 (10 Jul 2022)-img/p2_chem_s1_q05_D.png}{111.4}{0}
\qfigdef{p2_chem_s2_para9_10_1}{TOAS - Test 12 (10 Jul 2022)-img/p2_chem_s2_para9_10_1.png}{196.5}{0}
\qfigdef{p2_chem_s2_para9_10_A}{TOAS - Test 12 (10 Jul 2022)-img/p2_chem_s2_para9_10_A.png}{139.6}{0}
\qfigdef{p2_chem_s2_para9_10_B}{TOAS - Test 12 (10 Jul 2022)-img/p2_chem_s2_para9_10_B.png}{150.4}{0}
\qfigdef{p2_chem_s2_para9_10_C}{TOAS - Test 12 (10 Jul 2022)-img/p2_chem_s2_para9_10_C.png}{150.4}{0}
\qfigdef{p2_chem_s2_q09_A}{TOAS - Test 12 (10 Jul 2022)-img/p2_chem_s2_q09_A.png}{46.1}{0}
\qfigdef{p2_chem_s2_q09_B}{TOAS - Test 12 (10 Jul 2022)-img/p2_chem_s2_q09_B.png}{48.9}{0}
\qfigdef{p2_chem_s2_q09_C}{TOAS - Test 12 (10 Jul 2022)-img/p2_chem_s2_q09_C.png}{52.3}{0}
\qfigdef{p2_chem_s2_q09_D}{TOAS - Test 12 (10 Jul 2022)-img/p2_chem_s2_q09_D.png}{51.5}{0}
\qfigdef{p2_chem_s2_q10_1}{TOAS - Test 12 (10 Jul 2022)-img/p2_chem_s2_q10_1.png}{51.9}{0}
\qfigdef{p2_chem_s2_q10_2}{TOAS - Test 12 (10 Jul 2022)-img/p2_chem_s2_q10_2.png}{65.6}{0}
\qfigdef{p2_chem_s2_q10_4}{TOAS - Test 12 (10 Jul 2022)-img/p2_chem_s2_q10_4.png}{41.9}{0}
\qfigdef{p2_chem_s2_q10_3}{TOAS - Test 12 (10 Jul 2022)-img/p2_chem_s2_q10_3.png}{46.6}{0}
\qfigdef{p2_chem_s3_q03_1}{TOAS - Test 12 (10 Jul 2022)-img/p2_chem_s3_q03_1.png}{119.3}{0}
\qfigdef{p2_chem_s3_q03_2}{TOAS - Test 12 (10 Jul 2022)-img/p2_chem_s3_q03_2.png}{237.5}{0}
\qfigdef{p2_chem_s3_q03_3}{TOAS - Test 12 (10 Jul 2022)-img/p2_chem_s3_q03_3.png}{156.3}{0}
\qfigdef{p2_chem_s3_q03_4}{TOAS - Test 12 (10 Jul 2022)-img/p2_chem_s3_q03_4.png}{227.5}{0}
\raggedcolumns
\makeatletter\@beginparpenalty=10000\makeatother
\newcolumntype{L}[1]{>{\raggedright\arraybackslash}p{#1}}
\begin{document}
\begin{center}{\Large\bfseries TOAS --- Test 12 (10 Jul 2022)}\end{center}
\begin{multicols}{2}
\raggedcolumns
\papersec{1}
\partsec{1}{PHYSICS}
\sect{SECTION-I}
\qn{1} A charge particle $q_0$ is kept at centre of 2 non conducting hemi spherical shell of radius $R_1$ \& $R_2$. If $E_1$ \& $E_2$ is electric field at surfaces of A \& B then
\qfig{p1_phy_s1_q01_1}
\begin{qoptions}
\opt{A}{Ratio of flux $\dfrac{\phi_1}{\phi_2}=1$.}
\opt{B}{$E_1=E_2$}
\opt{C}{$\dfrac{R_1}{R_2}=\sqrt{\dfrac{E_2}{E_1}}$}
\opt{D}{$E_1\ne E_2$}
\end{qoptions}
\qn{2} A ray of light tangentialy incident on outer sphere.
\qfig{p1_phy_s1_q02_1}
\begin{qoptions}
\opt{A}{Net deviation suffered by ray is $240^\circ$.}
\opt{B}{Ray will never enter inside inner sphere.}
\opt{C}{Ray is incident tangentially on inner sphere.}
\opt{D}{Angle made by ray with normal inside inner sphere is $30^\circ$.}
\end{qoptions}
\qn{3} Variation of magnification of convex lens vs image distance for real objects is shown. Choose correct statement.
\qfig{p1_phy_s1_q03_1}
\begin{qoptions}
\opt{A}{$f$ is equal to intercept of V-axis.}
\opt{B}{$f$ is equal to $\dfrac{1}{\text{slope}}$ of graph.}
\opt{C}{$f$ is equal to slope of graph}
\opt{D}{magnitude of intercept on $m$ axis is 1.}
\end{qoptions}
\qn{4} Diagram shows a capacitor connected with a constant current source. Choose the correct option(s) :
\qfig{p1_phy_s1_q04_1}
\begin{qoptions}
\opt{A}{the rate of change of voltage across the capacitor in the circuit of figure from top to bottom 166.7 V/s.}
\opt{B}{the rate of change of voltage across the capacitor in the circuit of figure from top to bottom $-$166.7 V/s.}
\opt{C}{the capacitor's voltage 10 milliseconds after power is switched on is 1.667V.}
\opt{D}{the capacitor's voltage 10 milliseconds after power is switched on is 2.667V.}
\end{qoptions}
\qn{5} A long metallic pipe with internal diameter $d$ is pointing directly downwards. Water is slowly dripping from a nozzle at its lower end, see fig. Water can be considered to be electrically conducting; its surface tension is $\sigma$ and its density is $\rho$. A droplet of radius $r$ hangs below the nozzle. The radius grows slowly in time until the droplet separates from the nozzle due to the free fall acceleration $g$. Always assume that $d\ll r$.\par
Choose the correct option(s) :-
\qfig{p1_phy_s1_q05_1}
\begin{qoptions}
\opt{A}{the radius $r_{\max}$ of a drop just before it separates from the nozzle $\sqrt[3]{\dfrac14\dfrac{\sigma d}{\rho g}}$.}
\opt{B}{the radius $r_{\max}$ of a drop just before it separates from the nozzle $\sqrt[3]{\dfrac34\dfrac{\sigma d}{\rho g}}$.}
\opt{C}{Relative to the far-away surroundings, the pipe's electrostatic potential is $\phi$. the charge $Q$ of a drop when its radius is $r$, $Q=4\pi\varepsilon_0\phi r$}
\opt{D}{Consider the situation in which $r$ is kept constant and $\varphi$ is slowly increased. The droplet becomes unstable and breaks into pieces if the hydrostatic pressure inside the droplet becomes smaller than the atmospheric pressure, the critical potential $\varphi_{\max}$ at which this will happen $2\sqrt{\dfrac{\sigma r}{\epsilon_0}}$.}
\end{qoptions}
\qn{6} A capacitor consists of two circular parallel plates both with radius $R$ separated by distance $d$, where $d\ll R$, as shown in figure (a). The top plate is connected to a constant voltage source at a potential $V$ while the bottom plate is grounded. Then a thin and small disk of mass $m$ with radius $r$ ($\ll R,d$) and thickness $t$ ($\ll r$) is placed on the center of the bottom plate, as shown in figure (b).\par
Let us assume that the space between the plates is in vacuum with the dielectric constant $\epsilon_0$; the plates and the disk are made of perfect conductors; and all the electrostatic edge effects may be neglected. The inductance of the whole circuit and the relativistic effects can be safely disregarded. The image charge effect can also be neglected. Choose the correct option(s)
\qfig{p1_phy_s1_q06_1}
\begin{qoptions}
\opt{A}{When the disk is placed on the bottom plate, a charge $q$ on the disk of figure (b) is related to the voltage $V$ by $q=\chi V$, then $\chi$ in terms of $r$, $d$ and $\varepsilon_0$ is given by $-\dfrac{\epsilon_0\pi r^2}{d}$.}
\opt{B}{When the disk is placed on the bottom plate, a charge $q$ on the disk of figure (b) is related to the voltage $V$ by $q=\chi V$, then $\chi$ in terms of $r$, $d$ and $\varepsilon_0$ is given by $\dfrac{\epsilon_0\pi r^2}{d}$.}
\opt{C}{The parallel plates lie perpendicular to a uniform gravitational field $g$. To lift up the disk at rest initially, we need to increase the applied voltage beyond a threshould voltage $v_{th}$. The value of $v_{th}$ is given by $\sqrt{\dfrac{2mgd}{\chi}}$.}
\opt{D}{The parallel plates lie perpendicular to a uniform gravitational field $g$. To lift up the disk at rest initially, we need to increase the applied voltage beyond a threshould voltage $v_{th}$. The value of $v_{th}$ is given by $\sqrt{\dfrac{mgd}{\chi}}$.}
\end{qoptions}
\qn{7} A massive star of mass $M$ is in uniform circular orbit around a supermassive black hole of mass $M_b$. Initially, the radius and angular frequency of the orbit are $R$ and $\omega$ respectively. According to Einstein's theory of general relativity the space around the two objects is distorted and gravitational waves are radiated. Energy is lost through this radiation and as a result the orbit of the star shrinks gradually. One may assume, however, that the orbit remains circular throughout and Newtonian mechanics holds. Choose the correct option(s) :-
\begin{qoptions}
\opt{A}{The power radiated through gravitational wave by this star is given by $L_G=KC^xG^yM^2R^4\omega^6$ where $c$ is the speed of light, $G$ is the universal gravitational constant, and $K$ is a dimensionless constant, then $x=-5$ \& $y=1$}
\opt{B}{Total mechanical energy of the star-black hole system is $\dfrac{GMm_b}{2R}$}
\opt{C}{An expression for the rate of decrease in the orbital period ($dT/dt$) in terms of the masses, period $T$ and constants given by $\dfrac{dT}{dt}=-\dfrac{3KG^{5/3}M_b^{2/3}M(2\pi)^{8/3}}{c^5T^{5/3}}$}
\opt{D}{An expression for the rate of decrease in the orbital period ($dT/dt$) in terms of the masses, period $T$ and constants given by $\dfrac{dT}{dt}=\dfrac{3KG^{5/3}M_b^{2/3}M}{c^5T^{5/3}}$}
\end{qoptions}
\qn{8} An uncharged solid spherical conductor of radius $R$ has its centre at the origin of a cartesian frame OXY Z. Two point charges each of $+Q$ are placed at fixed positions at ($5R/2$, 0, 0) and (0, 0, $3R/2$). There will be induced charges on the sphere due to these charges. Choose the correct option(s)
\qfig{p1_phy_s1_q08_1}
\begin{qoptions}
\opt{A}{Electric field at $\left(\dfrac R2,\dfrac R2,\dfrac R2\right)$ due to induced charges is given by\par $E_p=-\dfrac{2^{3/2}Q}{(4\pi\epsilon_0R^2)}$ $\left[\left(\dfrac92-2\sqrt3\right)\hat i\right.$ $\left.+\left(\dfrac92+\sqrt3\right)\hat j\right.$ $\left.+(9-\sqrt3)\hat k\right]$}
\opt{B}{Electric field at $\left(\dfrac R2,\dfrac R2,\dfrac R2\right)$ due to induced charges is given by\par $E_p=-\dfrac{2^{3/2}Q}{(4\pi\epsilon_0R^2)}$ $\left[\left(\dfrac92-2\sqrt3\right)\hat i\right.$ $\left.+\left(\dfrac92+\sqrt3\right)\hat j\right.$ $\left.+(9+\sqrt3)\hat k\right]$}
\opt{C}{The potential at P due to induced charges is given by\par $V=\dfrac{Q}{(4\pi\epsilon_0R)\left(\frac{16}{15}-\left(\frac29\right)^{1/2}-\left(\frac23\right)^{1/2}\right)}$}
\opt{D}{The potential at P due to induced charges is given by\par $V=\dfrac{Q}{(4\pi\epsilon_0R)\left(\frac{16}{15}+\left(\frac29\right)^{1/2}+\left(\frac23\right)^{1/2}\right)}$}
\end{qoptions}
\qn{9} The graph shows variation of flux ($\phi$) with time ($t$) passing through a sphere for varries conditions. Choose the correct option(s) : (Assume : linear charge on x-axis)
\qfig{p1_phy_s1_q09_1}
\begin{qoptions}
\opt{A}{Graph `a' may corresponds to the case when a sphere is moving along x-axis with constant speed and the charge distribution is uniform in x-direction.}
\opt{B}{Graph `b' may corresponds to the case when a sphere is moving along x-axis with constant speed and the charge distribution is proportional to $x$.}
\opt{C}{Graph `c' may corresponds to the case when a sphere is moving along x-axis with constant speed and the charge distribution is varying linearly with $x$.}
\opt{D}{Graph `d' may corresponds to the case when a sphere is located at origin and its radius is increasing with uniform rate and charge distribution is proportional to $x$.}
\end{qoptions}
\qn{10} A non conducting rod AB of length $\sqrt3R$, uniformly distributed charge of linear charge density $\lambda$ and a non-conducting ring of uniformly distributed charge $Q$, are placed as shown in the figure. Point A is the centre of ring and line AB is the axis of ring, perpendicular to plane of ring. The electrostatic interaction energy between ring and rod is:
\qfig{p1_phy_s1_q10_1}
\begin{qoptions}
\opt{A}{$\dfrac{Q\lambda}{4\pi\varepsilon_0}\ln(2+\sqrt3)$}
\opt{B}{$\dfrac{Q\lambda}{2\pi\varepsilon_0}\ln(2+\sqrt3)$}
\opt{C}{$\dfrac{Q\lambda}{4\pi\varepsilon_0}\ln(2-\sqrt3)$}
\opt{D}{$\dfrac{Q\lambda}{2\pi\varepsilon_0}\ln(2-\sqrt3)$}
\end{qoptions}
\sect{SECTION-II}
\qn{1} In given circuit when switch S is in position A current in resistance is $i_1$ \& position B it is $i_2$, then $\dfrac{i_1}{i_2}$ is equal to:
\qfig{p1_phy_s2_q01_1}
\qn{2} x-z plane seprates 2 medium A \& B having $\mu_1=1.5$ \& $\mu_2=2$. A ray of light passes through two medium from A to B. If unit vector of ray are $a\hat i+b\hat j$ \& $c\hat i+d\hat j$ in A \& B. then value of $a/c$ is
\qn{3} A parallel plate capacitor consisting two dielectric materials is shown in the figure. The middle dielectric slab is place symmetrically with respect to the plates. If the potential difference between one of the plates and the nearest surface of dielectric interface is 2Volts, then the ratio $\epsilon_1:\epsilon_2$ is
\qfig{p1_phy_s2_q03_1}
\qn{4} Two ammeters X and Y have resistances of 1.2$\Omega$ and 1.5$\Omega$ respectively and they give full scale deflection with 150 mA and 250 mA respectively. The ranges have been extended by connecting shunts so as to give full scale deflection with 15 A. The ammeters along with shunts are connected in parallel and then placed in a circuit in which the total current flowing is 15A. The current in amperes indicated in ammeter X is.
\qn{5} Figure shows a uniformly charged disc of radius $\sqrt{39}$ m and charge density 3 $\mu$C/m$^2$ rotating about its axis with angular velocity 5 rad/sec. Find out the magnetic field at a distance 5m from the center of disc on its axis. Express your answer in pT.
\qfig{p1_phy_s2_q05_1}
\qn{6} Figure shows a ring of mass `M' and radius `R' uniformly charged with charge $Q$. A point charge `$-q$' of mass $m$ is kept at the centre of ring. Ring is free to move. Find the time period of small oscillation of the point charge. [Take : $m=1$ kg, $M=2$kg, $q=-3\mu$C, $Q=40\ \mu$C, $R=2$m]
\qfig{p1_phy_s2_q06_1}
\qn{7} A analog voltmeter uses external multiplier settings (an external resistance connected in series with voltmeter). With a multiplier setting of 20 k$\Omega$, it reads 440 V and with a multiplier setting of 80 k$\Omega$ it reads 352 V. For a multiplier setting of 40 k$\Omega$, find the reading (in V) of voltmeter
\qn{8} Helium is used as working substance in an engine working on the cycle as shown in the figure. Processes AB, BC, CD and DA are isobaric, adiabatic, isochoric and isothermal respectively. The ratio of maximum to minimum volume of helium during the cycle is $8\sqrt2$ and that of maximum to minimum absolute temperature is 4. Calculate efficiency ($\eta$) of the cycle ($\eta<1$).
\qfig{p1_phy_s2_q08_1}
\raggedcolumns
\partsec{2}{CHEMISTRY}
\sect{SECTION-I}
\qn{1} Which of the following statement(s) is/are \textbf{CORRECT}?
\begin{qoptions}
\opt{A}{In the bessemerization process of copper extraction from copper pyrite, ferrous silicate is formed in the form of slag}
\opt{B}{In hydrometallurgical process used for the extraction of copper from low grade ore, Fe is commercially used as reducing agent in displacement reaction but not Zinc}
\opt{C}{Se and Te is present in the anode mud during electrorefining of impure copper}
\opt{D}{Partial roasting of chalcopyrite produces Cu$_2$O + FeS}
\end{qoptions}
\qn{2} Select the incorrect statements (s) -
\begin{qoptions}
\opt{A}{Combustion of methane in an adiabatic rigid container will cause no change in temperature of system.}
\opt{B}{It is possible to have both adiabatic reversible and adiabatic irreversible process between two fixed states.}
\opt{C}{For reaction involving only ideal gas, occuring at constant temperature there will be no change in internal energy.}
\opt{D}{No process is possible for the gases in which volume is changed but PV work is zero.}
\end{qoptions}
\qn{3} In an evacuated vessel of 8.2 L capacity, 0.2 moles of A(g) is taken, vessel is sealed and heated to 127$^\circ$C, by which the following first order reaction occurred at constant volume and temperature :-\par
A(g) $\longrightarrow$ 2B(g) + C(s)\par
The vapour presure of C(s) is 0.1 atm at 127$^\circ$C. If the half life of reaction is 20 min, then the correct information(s) related with reaction is/are : ($R=0.082$ L-atm/K-mol)
\begin{qoptions}
\opt{A}{$P_{\text{total}}=1.7$ atm\quad at $t=\infty$}
\opt{B}{$P_{\text{total}}=0.98$ atm\quad after 10\% decomposition of A(g)}
\opt{C}{$P_{\text{total}}=0.88$ atm\quad after 5\% decompositon of A(g)}
\opt{D}{$P_{\text{total}}=1.3$ atm\quad at $t=20$ min}
\end{qoptions}
\qn{4} Which of the following relationship is(are) correct for the following figure.
\qfig{p1_chem_s1_q04_1}
\begin{qoptions}
\opt{A}{$x=y=z$}
\opt{B}{$x=y>z$}
\opt{C}{$z>120^\circ$}
\opt{D}{$x>y$}
\end{qoptions}
\qn{5} \qfig{p1_chem_s1_q05_A}
A and B are :-
\begin{qoptions}
\opt{A}{\qfigi{p1_chem_s1_q05_A2}}
\opt{B}{\qfigi{p1_chem_s1_q05_B}}
\opt{C}{\qfigi{p1_chem_s1_q05_C}}
\opt{D}{No formation of A and B}
\end{qoptions}
\qn{6} \qfig{p1_chem_s1_q06_1}
Products formed by P \& Q can not be differentiated by:
\begin{qoptions}
\opt{A}{2, 4 DNP}
\opt{B}{Lucas reagent (ZnCl$_2$) conc. HCl}
\opt{C}{NaHSO$_3$}
\opt{D}{Fehlings solution}
\end{qoptions}
\qn{7} Which of the following reaction(s) produce aldehyde as a product :
\begin{qoptions}
\opt{A}{\qfigi{p1_chem_s1_q07_A}}
\opt{B}{\qfigi{p1_chem_s1_q07_B}}
\opt{C}{CH$_3-$CN $\xrightarrow[\text{(ii) H}_3\text{O}^+]{\text{(i) SnCl}_2\text{ / HCl}}$}
\opt{D}{CH$_3$MgBr + CH$_3-$C$\equiv$N $\xrightarrow{\text{H}^+\text{ / H}_2\text{O}}$}
\end{qoptions}
\qn{8} Which of the following compound can only form
\qfig{p1_chem_s1_q08_1}
\begin{qoptions}
\opt{A}{\qfigi{p1_chem_s1_q08_A}}
\opt{B}{\qfigi{p1_chem_s1_q08_B}}
\opt{C}{\qfigi{p1_chem_s1_q08_C}}
\opt{D}{\qfigi{p1_chem_s1_q08_D}}
\end{qoptions}
\qn{9} Which of the following statement(s) is/are \textbf{INCORRECT}?
\begin{qoptions}
\opt{A}{H$_2$S when passed in sodium nitroprussiden solution then purple complex is formed}
\opt{B}{H$_2$S reacts with H$_2$SO$_3$ solution and produce sulphur , in this reaction H$_2$S acts as oxidizing agent}
\opt{C}{NaNO$_2$ produce a basic gas on reaction with Zn dust + KOH solution}
\opt{D}{NaNO$_2$ give red coloured product when treated with sulphanilic acid and dil. CH$_3$COOH followed by the addition of 1-nepthylamine.}
\end{qoptions}
\qn{10} Light from a discharge tube containing hydrogen atom fall on the surface of a piece of sodium. The kinetic energy of a fastest photo electron emitted from sodium is 0.73 eV and work function for sodium is 1.82 eV. Collector plate is connected to $+$ve terminal of battery having voltage 2 volts. Which of the following are true regarding above experiment :
\begin{qoptions}
\opt{A}{The energy of photon causing the photoelectric effect is 2.55 eV.}
\opt{B}{Hydrogen atoms electron undergoes a transition from $n=4$ to $n=2$ to show given photoelectric effect.}
\opt{C}{Change in angular momentum of the electron in the hydrogen atom will be $h/\pi$.}
\opt{D}{Kinetic energy of photoelectron at emitter plate will be 0.73 eV and at collector plate will be 2.73 eV.}
\end{qoptions}
\sect{SECTION-II}
\qn{1} If the number of following chemical species which form N$_2$O with dilute HNO$_3$ is $x$. Then predict the values of $x$.\par
Ag , Au , Fe$^{+2}$ , Zn , S , BaSO$_3$ , Hg , CaCO$_3$ , I$_2$
\qn{2} If the number of metal-metal bond, number of terminal CO ligand and number of bridging ligand in Co$_2$(CO)$_8$ (solid) are $x$, $y$ \& $z$ respectively.\par
Find $(x+y+z)$.
\qn{3} A fish swims in water at 300 K. In the process fish released 3000 Joule energy to water.\par
If $\Delta S_{\text{fish}}=-\dfrac23\times10^{-3}$J/K for every joule of energy released by fish, find $\Delta S_{\text{universe}}$ in the process (in \textbf{J/K})
\qn{4} Determine degree of dissociation (in \%) of 0.05 M NH$_3$ at 25$^\circ$C in a solution of pH = 11.
\qn{5} The vapour pressures of two pure liquids A and B which form an ideal solution are 300 and 800 Torr, respectively, at temperature $T$. A liquid solution of A and B for which the mole fraction of A is 0.60 is contained in a cylinder closed by a piston on which the pressure can be varied. The solution is slowly vapourized at temperature $T$ by decreasing the applied pressure. Find total vapour pressure (in torr) when last drop of liquid is remaining ?
\qn{6} Find sum of the oxidation number of Iodine in (X) and (Z).\par
KI + conc. H$_2$SO$_4\to$ (X) iodine containing product + other product\par
(X) $\xrightarrow[\text{hot/conc.}]{\text{KOH}}$ (Y) + (Z)\par
(Y) $\xrightarrow[\text{KMnO}_4]{\text{KOH}}$ (Z) + other product
\qn{7} \qfig{p1_chem_s2_q07_1}
How many ($x$) moles of HI consumed in the given reaction ?
\qn{8} How many compounds will give positive Fehling test :\par
(i) CH$_3$CHO\quad (ii) PhCHO\quad (iii) CH$_3-$CO$-$CH$_3$\quad (iv) CH$_3-$CH(OH)(OCH$_3$)
\qfig{p1_chem_s2_q08_2}
\raggedcolumns
\partsec{3}{MATHEMATICS}
\sect{SECTION-I}
\qn{1} Identify the correct statement(s) (where $C$ is the constant of integration).
\begin{qoptions}
\opt{A}{$\displaystyle\int\frac{(x^5-x)^{1/5}}{x^6}\,dx=\frac{5}{24}\left(1-\frac{1}{x^4}\right)^{6/5}+C$}
\opt{B}{$\displaystyle\int\frac{(x^5-x)^{1/5}}{x^6}\,dx=-\frac{5}{24}\left(1-\frac{1}{x^4}\right)^{6/5}+C$}
\opt{C}{$\displaystyle\int(\tan^3x+\tan^2x+\tan x+1)\,dx=\frac{(\tan x+1)^2}{2}+C$}
\opt{D}{$\displaystyle\int(\tan^3x+\tan^2x+\tan x+1)\,dx=\frac12\sec^2x+\tan x+C$}
\end{qoptions}
\qn{2} Given $I_n=\displaystyle\int_0^1(1-x^2)^n\,dx$ $(n=0,1,2,\dots)$, then which of the following is/are true?
\begin{qoptions}
\opt{A}{$\dfrac{I_{n+1}}{I_n}$ is independent of $n$}
\opt{B}{$\dfrac{I_{n+1}}{I_n}=\dfrac{2n+2}{2n+3}$}
\opt{C}{$I_k=\displaystyle\prod_{r=1}^{k}\left(\frac{2r}{2r+1}\right)$, $(k\in\mathbb{N})$}
\opt{D}{there exist some $n$ $(n\in\mathbb{N})$ for which $I_n\ge1$}
\end{qoptions}
\qn{3} $\displaystyle\int\sqrt{1+\operatorname{cosec}x}\;dx$ equals - (where $C$ is integration constant)
\begin{qoptions}
\opt{A}{$2\sin^{-1}\sqrt{\sin x}+C$}
\opt{B}{$\sqrt2\cos^{-1}\sqrt{\cos x}+C$}
\opt{C}{$C-2\sin^{-1}(1-2\sin x)$}
\opt{D}{$\cos^{-1}(1-2\sin x)+C$}
\end{qoptions}
\qn{4} Let $A_n=\displaystyle\sum_{k=0}^{n-1}\frac{n}{k^2+n^2+(n+2)k+n+1}$ and $B_n=\displaystyle\sum_{k=0}^{n-1}\frac{n}{n^2+kn+k^2}$ for $n=1,2,3,\dots$, then which of the following is/are true?
\begin{qoptions}
\opt{A}{$A_n<B_n$}
\opt{B}{$A_n>B_n$}
\opt{C}{$A_n<\dfrac{\pi}{3}$}
\opt{D}{$B_n>\dfrac{\pi}{9}$}
\end{qoptions}
\qn{5} If $f(x)=\displaystyle\lim_{n\to\infty}\frac{2x^{2n}+x+2}{x^{2n}-x^2+3}$, then which of the following statement(s) holds good?
\begin{qoptions}
\opt{A}{$f(x)$ is continuous at $x=1$}
\opt{B}{$f(x)$ is discontinuous at $x=1$}
\opt{C}{$\displaystyle\lim_{x\to2}f(x)$ exists finitely}
\opt{D}{$\displaystyle\lim_{x\to2}f(x)$ does not exist}
\end{qoptions}
\qn{6} The function $f(x)=[\cos x]\cos x$ in $(0,2\pi)$ -
\begin{qoptions}
\opt{A}{is discontinuous at two points}
\opt{B}{is non differentiable at two points}
\opt{C}{has range $(0,1]$}
\opt{D}{is such that $f(f(x))+1=0$ has no real roots}
\end{qoptions}
\qn{7} Let line $L$ is normal to parabola $y^2=4x$ and line $L$ intersect the parabola at $A$ and $B$. If point $P(\alpha,\beta)$ is point of intersection of tangents at $A$ and $B$ then which of the following is/are correct?
\begin{qoptions}
\opt{A}{Locus of point $P$ is symmetric w.r.t. $x$ axis}
\opt{B}{Locus of point $P$ is symmetric w.r.t. $y$ axis}
\opt{C}{If $\alpha,\beta\in I$ then total number of possible ordered pair $(\alpha,\beta)$ is 2}
\opt{D}{If $\alpha,\beta\in I$ then total number of possible ordered pair $(\alpha,\beta)$ is 4}
\end{qoptions}
\qn{8} Let $f(x)=ax^9+b\sin x+cx^2\operatorname{sgn}(x)+\dfrac{e^x-e^{-x}}{e^x+e^{-x}}$ be defined on set of real number (where $a>0$, $b,c\in\mathbb{R}$). If $f(-5)=5$, $f(-2)=-3$, then the minimum number of zeroes of the equation $f(x)=0$, is
\begin{qoptions}
\opt{A}{1}
\opt{B}{2}
\opt{C}{3}
\opt{D}{5}
\end{qoptions}
\qn{9} Given $f$ and $g$ are twice differentiable functions and $f'\neq0$, $g'\neq0$. If $f(x)g(x)=1$, then $\dfrac{f''(x)}{f'(x)}-\dfrac{g''(x)}{g'(x)}$ is equal to :
\begin{qoptions}
\opt{A}{$\dfrac{-2g'(x)}{g(x)}$}
\opt{B}{$\dfrac{g'(x)}{g(x)}$}
\opt{C}{$\dfrac{2f'(x)}{f(x)}$}
\opt{D}{$\dfrac{-f'(x)}{f(x)}$}
\end{qoptions}
\qn{10} If $\displaystyle\int\frac{7+6x^{13}}{\sqrt[3]{x^{17}-x^{30}}}\,dx=-\dfrac{p\,(1-x^{13})^{q/3}}{2x^{r/3}}+C$ where $p,q,r\in\mathbb{N}$ and $C$ is the constant of integration, then
\begin{qoptions}
\opt{A}{$p=3$}
\opt{B}{$p=7$}
\opt{C}{$r-q=12$}
\opt{D}{$r-q=8$}
\end{qoptions}
\sect{SECTION-II}
\qn{1} If $\displaystyle\int_0^1\frac{e^x\,dx}{\sqrt{1-x^2}}=A$ and $\displaystyle\int_0^\pi\left(e^{|\sin x|}+e^{|\cos x|}\right)dx$ is equal to $\lambda A$, then the value of $\lambda$ is -
\qn{2} The value of $\displaystyle\int_0^1\sqrt[3]{2x^3-3x^2-x+1}\;dx$ is -
\qn{3} If $\displaystyle\int\frac{(3x^2+2x)\,dx}{x^6+2x^5+x^4+2x^3+2x^2+5}=\frac12\tan^{-1}\bigl(g(x)\bigr)+C$, then the value of $\bigl(g(1)+g(0)\bigr)$ is
\qn{4} Let $f(x)$ be a twice differentiable function such that $f''(0)=2$, then the value of $\displaystyle\lim_{x\to0}\frac{f(5x)-3f(3x)+2f(2x)}{x^2}$ is
\qn{5} Let $g:\left[\tfrac12,1\right]\to(0,\infty)$, $g(x)=e^{-2x}f(x)$, where $g(x)$ is non constant decreasing and differentiable function and $g\!\left(\tfrac12\right)=\dfrac1e$. Then the value of $\left[\displaystyle\int_{1/2}^{1}f(x)\,dx\right]$ is\par
[where $[x]$ denotes the greatest integer less than or equal to real number $x$]
\qn{6} The number of values of $x$, $x\in[-2,3]$ where $f(x)=[x^2]\sin(\pi x)$ is discontinuous, is\par
(where $[\,.\,]$ denotes greatest integer function) -
\qn{7} Let $\tan^{-1}x=\dfrac12\tan^{-1}y=\dfrac13\tan^{-1}z=\dfrac{\pi}{7}$, then the value of $\dfrac12\left(x^2+y^2+z^2-x^2y^2z^2\right)$, is
\qn{8} The maximum value of $f(x)=\dfrac{2x^2+7x+3}{2x^2+x+1}$, is
\raggedcolumns
\newpage
\papersec{2}
\partsec{1}{PHYSICS}
\sect{SECTION-I(i)}
\qn{1} One mole of an ideal gas is expanded from the state A $(P_0,V_0)$ to final state B having volume $V$. The process follows a path represented by a straight line on the P-V diagram as shown. Up to what volume $(V)$ the gas shall be expanded so that final internal energy is half the maximum internal energy during the process.
\qfig{p2_phy_s1_q01_1}
\begin{qoptions}
\opt{A}{$\left(\dfrac{\sqrt2+1}{\sqrt2}\right)V_0$}
\opt{B}{$\sqrt2V_0$}
\opt{C}{$(\sqrt2+1)V_0$}
\opt{D}{$\left(\dfrac{\sqrt2+1}{2}\right)V_0$}
\end{qoptions}
\qn{2} A straight segment OP of length $L$ of a circuit carrying current $I$ ampere is placed along x-axis. Two infinitely long straight wires A and B each extending from $z=-\infty$ to $+\infty$ are fixed at $y=-\ell$ meter and $y=+\ell$ meter respectively. Wires A and B, each carry current $I$ ampere along positive z-axis. Given that O is origin of the coordinate system. The magnitude of force on segment OP is :-
\begin{qoptions}
\opt{A}{$\dfrac{\mu_0}{\pi}I^2\ln\left(1+\dfrac{L^2}{\ell^2}\right)$}
\opt{B}{$\dfrac{\mu_0}{2\pi}I^2\ln\left(1+\dfrac{L^2}{\ell^2}\right)$}
\opt{C}{$\dfrac{\mu_0}{\pi}I^2\ln\left(1+\dfrac{\ell^2}{L^2}\right)$}
\opt{D}{$\dfrac{\mu_0}{2\pi}I^2\ln\left(1+\dfrac{\ell^2}{L^2}\right)$}
\end{qoptions}
\qn{3} In the diagram shown, the charge $+Q$ is fixed. Another charge $+2q$, is projected from a distance $R$ from the fixed charge. The minimum separation between the two charges when the velocity becomes $\dfrac{1}{\sqrt3}$ times of the projected velocity at that instant, is (Assume gravity to the absent)
\qfig{p2_phy_s1_q03_1}
\begin{qoptions}
\opt{A}{$\dfrac{\sqrt3}{2}R$}
\opt{B}{$\sqrt3R$}
\opt{C}{$\dfrac12R$}
\opt{D}{$4R$}
\end{qoptions}
\qn{4} An equi-convex lens is made of a material of refractive index $\mu_1$ and is placed inside two liquids of refractive indices $\mu_2$ and $\mu_3$ as shown with $\mu_2>\mu_1>\mu_3$. A wide, parallel beam of light is incident on the lens from the left. The lens will give rise to :
\qfig{p2_phy_s1_q04_1}
\begin{qoptions}
\opt{A}{A single convergent beam.}
\opt{B}{Two different convergent beams.}
\opt{C}{Two different divergent beams}
\opt{D}{A convergent and divergent beam.}
\end{qoptions}
\qn{5} A uniformly charged non-conducting rod is suspended vertically at its end. The rod can swing freely in the vertical plane without any friction. The linear charge density on the rod is $\lambda$ C/m and it has a uniformly distributed mass $M$. Length of the rod is $L$. A uniform horizontal field $E$ is switched on so that the rod just manages to make itself horizontal. What is the maximum angular speed of the rod during its motion ?
\begin{qoptions}
\opt{A}{$\sqrt{\dfrac{3(\sqrt3-1)g}{L}}$}
\opt{B}{$\sqrt{\dfrac{3(\sqrt2-1)g}{2L}}$}
\opt{C}{$\sqrt{\dfrac{3(\sqrt2-1)g}{L}}$}
\opt{D}{$\sqrt{\dfrac{2(\sqrt3-1)g}{L}}$}
\end{qoptions}
\qn{6} In the circuit shown, the switch S is open initially. The capacitor $C$ is uncharged but $C_0$ has a charge $q_0$ at $t=0$
\qfig{p2_phy_s1_q06_1}
\begin{qoptions}
\opt{A}{The current immediately after switch S is closed is $I_C=\dfrac{\left(E+\frac{q_0}{C_0}\right)}{R}$.}
\opt{B}{The current immediately after switch S is closed is $I_c=\dfrac{E}{R}$}
\opt{C}{Final charge on $C$ after closing switch is $\dfrac{EC_0C}{C_0+C}$}
\opt{D}{Final charge on $C$ after closing switch is $\dfrac{EC_0C}{C_0+C}\left(1-\dfrac{q_0}{C_0E}\right)$}
\end{qoptions}
\qfullbegin
\sect{SECTION-I(ii)}
\paraq{7}{8}
Gaseous products of burning are released into the atmosphere of temperature $T_a$ through a high chimney of cross section $A$ and height $h$ as shown in the figure. The solid matter is burnt in the furnace which is at temperature $T_S$. due to burning smoke (gas) generation rate (volume per unit time) is $Q$. It can be assume that
\qfig{p2_phy_s2_para7_8_1}
\qfullsplit
$\to$ the velocity of the smoke (gas) in the furnace is negligibly small.\par
$\to$ the density of the smoke (gas) does not differ from that air at the same temperature and pressure.\par
$\to$ the gases can be treated as ideal while in furnace.\par
$\to$ the pressure of the air changes with height according to the hydrostatic law. The change of the density of the air with height is negligible.\par
$\to$ the flow of gases in the chimney follows Bernoulli's equation.\par
$\to$ the change of the density of the gas (smoke) is negligible throughout the chimney.
\qfullend
\qn{7} What is the minimum height of the chimney needed in order that chimney functions efficiently, so that it can release all of the produced smoke (gas) into the atmosphere ? Express your answer in terms of $Q$, $A$, $g$, $T_a$, $\rho_A$, $\rho_S$ and $\Delta T$, where $\Delta T=T_S-T_a$
\begin{qoptions}
\opt{A}{$\dfrac{\rho_ST_aQ^2}{gA^2\Delta T\,(\rho_a-\rho_S)}$}
\opt{B}{$\dfrac{T_aQ^2}{2gA^2\Delta T}$}
\opt{C}{$\dfrac{Q^2\rho_S}{2\,(\rho_a-\rho_S)A^2g}$}
\opt{D}{$\dfrac{Q^2\,(T_a+\Delta T)}{2gA^2\Delta T}$}
\end{qoptions}
\qn{8} Assume that two chimneys are built to serve exactly the same purpose. Their cross sections are identical, but are designed to work in different parts of the world, one in cold regions designed to work at an average atmospheric temperature of $-23^\circ$C and the other in warm regions, designed to work at an average atmospheric temperature of $27^\circ$C. The temperature of the furnace is $327^\circ$C. It was calculated that the height of the chimney designed to work in cold regions is 100 m. How high is the other chimney ?
\begin{qoptions}
\opt{A}{71.4 m}
\opt{B}{280 m}
\opt{C}{140 m}
\opt{D}{100 m}
\end{qoptions}
\paraq{9}{10}
A point object is placed at a distance 50 cm from the pole of a concave mirror. The radius of curvature of mirror is 30 cm. Point object O oscillates with amplitude of 1mm perpendicular to the principal axis as shown.
\qfig{p2_phy_s2_para9_10_1}
\qn{9} The amplitude of image is :-
\begin{qoptions}
\opt{A}{$\dfrac37$ mm}
\opt{B}{$\dfrac27$ mm}
\opt{C}{$\dfrac43$ mm}
\opt{D}{$\dfrac{11}{7}$ mm}
\end{qoptions}
\qn{10} Phase difference between motion of object and image when object crosses the principal axis is :-
\begin{qoptions}
\opt{A}{$\pi$}
\opt{B}{0}
\opt{C}{$\dfrac\pi2$}
\opt{D}{None of these}
\end{qoptions}
\paraq{11}{12}
Two moles of an ideal mono-atomic gas is taken through a cyclic process ABCA as shown in the V-T diagram. In the process BC, $VT=$ constant. The temperature at A is $T_0=300$ K. [take $\ln(3)=1.1$]. Then answer the following questions.
\qfig{p2_phy_s2_para11_12_1}
\qn{11} The work done by the gas during the process `BC' is :-
\begin{qoptions}
\opt{A}{600 R}
\opt{B}{900 R}
\opt{C}{1200 R}
\opt{D}{1500 R}
\end{qoptions}
\qn{12} The efficiency (in percentage) of the cycle is :-
\begin{qoptions}
\opt{A}{12}
\opt{B}{18}
\opt{C}{24}
\opt{D}{36}
\end{qoptions}
\sect{SECTION-II}
\qn{1} In the circuit shown in the figure, find the heat generated (in $\mu$J) through the resistor upto the instant when charge on capacitor A is half of the maximum values that the capacitors would store finally. (Initially both the capacitor are uncharged) (Given : $C=8\ \mu$F, $\varepsilon=3$ volt).
\qfig{p2_phy_s3_q01_1}
\qn{2} In a gravity free space, three identical metal plates have been placed parallel to each other as shown. In the initial state shown in the diagram, the potential difference between the points A and B is $V_1$. Now, the middle plate is shifted down a distance $L/2$, keeping the battery connected. The other plates are kept fixed. The potential difference between the points A and B now becomes $V_2$. The value $\dfrac{6V_1}{V_2}$ is equal to
\qfig{p2_phy_s3_q02_1}
\qn{3} If the optical axis of convex and concave lenses are separated by a distance 5 mm shown in the figure. If the coordinates of final image formed by the combination of Lenses is $(x,y)$ when the Parallel Beam of Light is incident on Convex Lens. If origin is at the Optical centre of convex lens, find $x+4y$ (in cm).
\qfig{p2_phy_s3_q03_1}
\qn{4} Find the equivalent resistance (in $\Omega$) between A and B in the given network.
\qfig{p2_phy_s3_q04_1}
\qn{5} A solid sphere of radius `R' is uniformly charged with charge density $\rho$ in its volume. A spherical cavity of radius $\dfrac R2$ is made in the sphere as shown in the figure. Electric potential at the centre of the sphere is $\dfrac{A\rho R^2}{B\epsilon_0}$. Where A and B are positive integers. Find $A^2B$.
\qfig{p2_phy_s3_q05_1}
\qn{6} Two particles A and B of mass 1kg and 2kg respectively are kept 1 meter apart and are released to move under mutual attraction. When speed of B is 3.6 cm/hour, what is the separation (in cm) between the particles. Take $G=\dfrac{20}{3}\times10^{-11}$ Nm$^2$/kg.
\raggedcolumns
\partsec{2}{CHEMISTRY}
\sect{SECTION-I(i)}
\qn{1} Consider the following reactions :-\par
A(g) $\xrightarrow{\text{FeCl}_3}$ (B)$\downarrow$ $+$ coloured ion\par
A(g) $\xrightarrow{\text{Cr}_2\text{O}_7^{2-}/\text{H}^{\oplus}}$ (B)$\downarrow$ $+$ coloured ion\par
C(aq.) $\xrightarrow{\text{Zn}+\text{dil. H}_2\text{SO}_4}$ A(g)\par
Identify A, B \& C respectively.
\begin{qoptions}
\opt{A}{SO$_2$, H$_2$S, S$_2$O$_3^{-2}$}
\opt{B}{H$_2$S, S$_8$, SO$_2$}
\opt{C}{SO$_2$, S$_8$, S$^{-2}$}
\opt{D}{SO$_2$, S$_8$, SO$_4^{-2}$}
\end{qoptions}
\qn{2} Consider the following reaction $2A(g)\to3B(g)+C(g)$. Starting with pure A having pressure 2 atm initially, the total pressure is exactly doubled in 2 hrs. The possible order of reaction is :-
\begin{qoptions}
\opt{A}{zero}
\opt{B}{first}
\opt{C}{second}
\opt{D}{third}
\end{qoptions}
\qn{3} The compressibility factor of a Vanderwaal gas is 0.5 at 27$^\circ$C and 24 atm. The value of Vanderwaal constant `a' (in atm-L$^2$/mol$^2$) is\par
($b=0.10$ L/mol, $R=0.08$ L-atm/K-mol)
\begin{qoptions}
\opt{A}{2.4}
\opt{B}{4.0}
\opt{C}{3.6}
\opt{D}{9.0}
\end{qoptions}
\qfullbegin
\qn{4}
\qfig{p2_chem_s1_q04_1}
\qfullsplit
R$_1$, R$_2$ \& R$_3$ are respectively :
\begin{qoptions}
\opt{A}{LiAlH$_4$, NaBH$_4$, B$_2$H$_6$}
\opt{B}{B$_2$H$_6$, NaBH$_4$, LiAlH$_4$}
\opt{C}{NaBH$_4$, LiAlH$_4$, B$_2$H$_6$}
\opt{D}{LiAlH$_4$, B$_2$H$_6$, NaBH$_4$}
\end{qoptions}
\qfullend
\qn{5} Which products are formed when compound \qfigi{p2_chem_s1_q05_1} is treated with concentrated aqueous KOH solution?
\begin{qoptions}
\opt{A}{\qfigi{p2_chem_s1_q05_A}}
\opt{B}{\qfigi{p2_chem_s1_q05_B}}
\opt{C}{\qfigi{p2_chem_s1_q05_C}}
\opt{D}{\qfigi{p2_chem_s1_q05_D}}
\end{qoptions}
\qn{6} Assuming that during the formation of XY$^{\oplus}$ from XY the bond order increase where as during the formation of Z$_2^{\oplus}$ from Z$_2$ the bond order decrease.\par
Based on above information which of the following statement is incorrect?
\begin{qoptions}
\opt{A}{During the formation of XY$^+$ from XY the electron is lost from antibonding molecular orbital}
\opt{B}{During the formation of Z$_2^{\oplus}$ from Z$_2$ the electron is lost from bonding molecular orbital}
\opt{C}{In above change XY undergoes in oxidation}
\opt{D}{In above change Z$_2$ undergoes in reduction}
\end{qoptions}
\sect{SECTION-I(ii)}
\paraq{7}{8}
The crystal field splitting, $\Delta_o$ depends upon the field produced by the ligand and charge on the metal ion. Some ligands are able to produce strong fields in which case, the splitting will be large whereas others produce weak field and consequently result in small splitting of d-orbitals.\par
The influence of the ligand on the colour of a complex may be, illustrated by considering the [Ni(H$_2$O)$_6$]$^{+2}$ complex. (which forms when nickel (II) chloride is dissolved in water) and reaction of it with ethane 1, 2-diamine (en).
\qn{7} If the bidentate ligand ethane 1,2-diamine (en) is progressively added in green aqueous solution of [Ni(H$_2$O)$_6$]Cl$_2$ in the molar ratios en : Ni as, 1 : 1 , 2 : 1 , 3 : 1 then their associated colour are respectively :
\begin{qoptions}
\opt{A}{Pale blue, Violet, Blue}
\opt{B}{Violet, Pale blue, Blue}
\opt{C}{Violet, Blue, Pale blue}
\opt{D}{Pale Blue, Blue, Violet}
\end{qoptions}
\qn{8} Sum of the stereo isomeric forms possible for complexes which are formed in above question :
\begin{qoptions}
\opt{A}{5}
\opt{B}{2}
\opt{C}{4}
\opt{D}{3}
\end{qoptions}
\paraq{9}{10}
Aldehydes and ketones react with an alcohol from compounds called hemiacetals, in which there is one hydroxyl group and one either-like group. Reaction of a hemiacetal with a second molecule of alchol gives an acetal and a molecule of water. We study this reaction
\qfig{p2_chem_s2_para9_10_1}
\noindent The carbonyl group of an aldehyde or ketone \quad A hemiacetal (has an --OH and an --OR group to the same carbon) \quad An acetal (has two --OR groups to the same carbon)\par
Draw structural formulas for the hemiacetal and acetal formed from these reagents. The stoichiometry of each reaction is given in the problem.\par
\noindent(A)
\qfig{p2_chem_s2_para9_10_A}
\noindent(B)
\qfig{p2_chem_s2_para9_10_B}
\noindent(C)
\qfig{p2_chem_s2_para9_10_C}
\qn{9} Product P is not possible -
\begin{qoptions}
\opt{A}{\qfigi{p2_chem_s2_q09_A}}
\opt{B}{\qfigi{p2_chem_s2_q09_B}}
\opt{C}{\qfigi{p2_chem_s2_q09_C}}
\opt{D}{\qfigi{p2_chem_s2_q09_D}}
\end{qoptions}
\qn{10} Which of the following statement is/are correct?
\begin{qoptions}
\opt{A}{Product Q is \qfigi{p2_chem_s2_q10_1}}
\opt{B}{Product Q is \qfigi{p2_chem_s2_q10_2}}
\opt{C}{Product R is \qfigi{p2_chem_s2_q10_4}}
\opt{D}{Product R is \qfigi{p2_chem_s2_q10_3}}
\end{qoptions}
\paraq{11}{12}
Enthalpy of combustion of a given compounds is the change in enthalpy when 1 mole of this compound combines with required amount of O$_2$ to give product in their stable forms. The enthalpy of combustion is usually measured by placing a known mass of the compound in a closed steel container (known as bomb-calorimeter). Graphite is thermodynamically most stable form of carbon.
\qn{11} Given the data below :\par
$\Delta H^0_{\text{comb}}$(CH$_4$, g) $=-900$ kJ/mole\par
$\Delta H^0_{\text{comb}}$(C, graphite) $=-400$ kJ/mole\par
$\Delta H^0_{\text{comb}}$(H$_2$, g) $=-300$ kJ/mole\par
$\Delta H^0_{\text{atomization}}$(C, graphite) $=700$ kJ/mole\par
and $\Delta H^0_f$(H, g) $=220$ kJ/mole\par
Identify the incorrect option(s) :
\begin{qoptions}
\opt{A}{Bond enthalpy of C--H bond is 340 kJ/mole}
\opt{B}{1 mole graphite produces more heat on combustion than 1 mole diamond at 1 atm and 25$^\circ$C}
\opt{C}{Enthalpy of atomisation of CH$_4$(g) is 1600 kJ/mole}
\opt{D}{Bond enthalpy of C--H bond is 420 kJ/mole}
\end{qoptions}
\qn{12} Identify the correct statement -
\begin{qoptions}
\opt{A}{$\Delta H^0_{\text{comb}}$(C, graphite) is equal to $\Delta H^0_f$(CO$_2$, g)}
\opt{B}{$\Delta H^0_f$(H$_2$O, $\ell$) is equal to $\Delta H^0_{\text{comb}}$(H$_2$, g)}
\opt{C}{The standard enthalpy of combustion is always negative.}
\opt{D}{The standard enthalpy of formation is always negative}
\end{qoptions}
\sect{SECTION-II}
\qn{1} A solution containing 30 g of a non volatile solute in exactly 90 g water has a vapour pressure of 20 mm Hg at 25$^\circ$C. Further 18 g of water is then added to the solution. The resulting solution has vapour pressure of 21 mm Hg at 25 $^\circ$C. Calculate molar mass of the solute.
\qn{2} If zinc blende (ZnS) crystal is heated, all ions present on any one face are lost. Monovalent anion (Q$^-$) enters each unit cell to maintain electrical neutrality without disturbance. If new formula of unit cell becomes Zn$_x$S$_y$Q$_Z$, then $(x+y+z)$ is :-
\qn{3} Identify number of reactions correctly match with their products ?\par
(a) PhCH$=$O $\xrightarrow[\text{D}_2\text{O}]{\text{NaOD}}$ PhCH$_2$OD $+$ PhCO$_2^-$
\qfig{p2_chem_s3_q03_1}
\qfig{p2_chem_s3_q03_2}
\qfig{p2_chem_s3_q03_3}
\qn{4}
\qfig{p2_chem_s3_q03_4}
(j) CH$_3$--CH(OH)--Ph\par
Answer in two digit form XY\par
X = Number of compounds giving (+) test with AgNO$_3$ / NH$_4$OH\par
Y = Number of compounds giving yellow colour with I$_2$+ NaOH
\qn{5} Consider the following compounds.\par
Pb(NO$_3$)$_2$, LiNO$_3$, NaNO$_3$, NH$_4$NO$_2$, KMnO$_4$, (NH$_4$)$_2$Cr$_2$O$_7$, HNO$_3$ (conc.) (NH$_4$)$_2$CO$_3$\par
If total number of compound(s) which produce coloured gas on heating $=x$\par
If total number of compound(s) which produce atleast one paramagnetic gaseous product $=y$\par
then find $\dfrac{x}{y}=$?
\qn{6} Consider the following anions?\par
CO$_3^{-2}$, SO$_3^{-2}$, S$^{-2}$, S$_2$O$_3^{-2}$, NO$_2^-$, HCO$_3^-$, HSO$_3^-$, CH$_3$COO$^-$\par
If total number of anion which produce colourless solution with KMnO$_4$/H$^{\oplus}$ by redox reaction $=A$\par
If total number of anion which produce coloured volatile product with dil. HCl $=B$\par
Then find $\dfrac{A+B}{2}=$?
\raggedcolumns
\partsec{3}{MATHEMATICS}
\sect{SECTION-I(i)}
\qn{1} If $\dfrac{\cos^4x}{\theta_1}+\dfrac{\sin^4x}{\theta_2}=\dfrac{1}{(\theta_1+\theta_2)}$, then $\dfrac{d\theta_2}{d\theta_1}$ equals (where $x$ is not dependent on $\theta_1$ or $\theta_2$)
\begin{qoptions}
\opt{A}{$\sin^2x$}
\opt{B}{$\cos^2x$}
\opt{C}{$\tan^2x$}
\opt{D}{$\sec^2x$}
\end{qoptions}
\qn{2} Let $F(x)=C\sin x\displaystyle\int_0^x\cos t\,dt+2\int_0^x t\,dt+C\cos^2x-x^2$ (where $C\in R^+$). If $x^2-2x+3\ge F(x)\ \forall\ x\in R$, then the maximum value of the function $\left(\dfrac{6e^x-e^{2x}-1}{e^x}\right)^{F(2)}$ is equal to -
\begin{qoptions}
\opt{A}{2}
\opt{B}{4}
\opt{C}{16}
\opt{D}{32}
\end{qoptions}
\qn{3} Suppose that $f$ is differentiable function such that $f'(x)\le2$ for all $x$. If $f(1)=2$ and $f(4)=8$ then $f(2)$ has the value equal to
\begin{qoptions}
\opt{A}{4}
\opt{B}{2}
\opt{C}{1}
\opt{D}{0}
\end{qoptions}
\qn{4} Let $f(x)$ be a differentiable function in $[-1,\infty)$ and $f(0)=1$ such that $\displaystyle\lim_{t\to x+1}\dfrac{t^2f(x+1)-(x+1)^2f(t)}{f(t)-f(x+1)}=1$, then the value of $\displaystyle\lim_{x\to1}\dfrac{\ln\,(f(x))-\ln2}{x-1}$ is
\begin{qoptions}
\opt{A}{0}
\opt{B}{1}
\opt{C}{2}
\opt{D}{does not exist}
\end{qoptions}
\qn{5} If $f(x)=(\sin x)^{1/3}\cdot(\sin2x)^{1/2}\cdot\sin3x$, then value of $f'\!\left(\dfrac\pi3\right)$ is of the form $-\dfrac{3^a}{2^b}$, then value of $\left(a-\dfrac b2\right)$ is
\begin{qoptions}
\opt{A}{1}
\opt{B}{2}
\opt{C}{5}
\opt{D}{0}
\end{qoptions}
\qn{6} Let $\displaystyle\int\dfrac{x^2+20}{(x\sin x+5\cos x)^2}\,dx=f(x)+\tan x+C$, $f(0)=0$ where $C$ denotes constant of integration. Then the value of $\pi+5f(\pi)$ is equal to
\begin{qoptions}
\opt{A}{2}
\opt{B}{4}
\opt{C}{3}
\opt{D}{0}
\end{qoptions}
\sect{SECTION-I(ii)}
\paraq{7}{8}
Let $f(x)=3x^2$ and $g(x)=2x^3+1$ are two curves and $L$ is the equation of tangent at the point where $f(x)$ and $g(x)$ touch each other
\qn{7} Which option(s) is/are correct -
\begin{qoptions}
\opt{A}{Equation of $L$ is $6x-y=3$}
\opt{B}{Equation of $L$ is $x+2y=7$}
\opt{C}{Equation of normal to the curve $g(x)$ from $(7,2)$ is $x+6y=19$}
\opt{D}{Equation of normal to the curve $f(x)$ from $(7,2)$ is $x+6y=19$}
\end{qoptions}
\qn{8} Which option(s) is/are correct -
\begin{qoptions}
\opt{A}{Angle between $f(x)$ and $g(x)$ at the point of intersection is $\tan^{-1}\dfrac97$}
\opt{B}{Angle between $f(x)$ and $g(x)$ at the point of intersection is $\sin^{-1}\dfrac{9}{\sqrt{130}}$}
\opt{C}{Point of intersection of $L$ and $g(x)$ is $(-2,-15)$}
\opt{D}{Value of $\displaystyle\lim_{x\to1}\dfrac{g(x)-f(x)}{(x-1)^2}$ is 3}
\end{qoptions}
\paraq{9}{10}
Let $f(x)=ax^3+bx^2+cx+d$, satisfying $\displaystyle\lim_{x\to0}\left(3+\dfrac{f(x)}{x^2}\right)^{\frac1x}=e^2$, then\par
On the basis of above information, answer the following questions :
\qn{9} Correct option's is /are -
\begin{qoptions}
\opt{A}{$a+b=0$}
\opt{B}{$a+b+c+d=0$}
\opt{C}{$a+c=2$}
\opt{D}{$d-b=2$}
\end{qoptions}
\qn{10} Sum of all roots of the all possible equations $f(x)-2(r^2-1)x^2=0$ (where $r=1,2,\dots,10$)
\begin{qoptions}
\opt{A}{385}
\opt{B}{250}
\opt{C}{165}
\opt{D}{305}
\end{qoptions}
\paraq{11}{12}
Consider $f:R\to R$ be a differentiable function such that $f(x)=x+\displaystyle\int_0^x e^{-t}f(x-t)\,dt$.
\qn{11} Minima of function $y=f(x)$ occurs at
\begin{qoptions}
\opt{A}{$x=-1$}
\opt{B}{$x=\dfrac12$}
\opt{C}{$x=1$}
\opt{D}{$x=2$}
\end{qoptions}
\qn{12} The value of $\displaystyle\int_1^2\dfrac{f''(x)}{x^2}dx$ is less than
\begin{qoptions}
\opt{A}{1}
\opt{B}{$\dfrac14$}
\opt{C}{$-\dfrac12$}
\opt{D}{$\dfrac34$}
\end{qoptions}
\sect{SECTION-II}
\qn{1} If $\displaystyle\int\dfrac{e^x\,(x^2-x+1)}{(x^2+1)\sqrt{x^2+1}}\,dx=e^x\cdot f(x)+C$, where $C$ is constant of integration, then number of integers in the range of $f(x)$ is
\qn{2} If $n\in N$ and $f(x)=({}^nC_0x+1)({}^nC_1x+2)({}^nC_2x+3)({}^nC_3x+4)\dots({}^nC_nx+n+1)$ and $f'(0)=\dfrac32(7!)$, then $n$ is equal to
\qn{3} Let $P(x)$ be a polynomial of degree 4 having a local maximum at $x=2$ and $\displaystyle\lim_{x\to0}\left(3-\dfrac{P(x)}{x}\right)=27$. If $P(1)=-9$ and $P''(x)$ has a local minimum at $x=2$. If global maximum value of $y=P'(x)$ on the set $A=\{x:x^2+12\le7x\}$ is $4M$ then value of $M$ is :
\qn{4} If $0<x<1$, then the solution of equation $3\sin^{-1}\!\left(\dfrac{2x}{1+x^2}\right)+2\tan^{-1}\!\left(\dfrac{2x}{1-x^2}\right)-4\cos^{-1}\!\left(\dfrac{1-x^2}{1+x^2}\right)=\dfrac\pi4$ is
\qn{5} For $f(x)=x^4+|x|$, Let $I_1=\displaystyle\int_0^\pi f(\cos x)\,dx$, $I_2=\int_0^{\pi/2}f(\sin x)\,dx$, then the value of $\dfrac{I_1}{I_2}$ is
\qn{6} If $f(x)=\begin{cases}x^2-5x+6, & 2<x<3\\ px^2+qx+1, & 3\le x\end{cases}$ is differentiable at $x=3$, then value of $(p-q)$ is
\end{multicols}
\end{document}
